% VI. Експериментални резултати и анализ.
% Всяко число тук идва от изпълнение на sentinel_bench на машината от Таблица 5.1.
% Пълният необобщен изход е в docs/measurements/bench_run.txt и в Приложение Г.
\section{Експериментални резултати и анализ}
\label{sec:results}

Разделът представя получените резултати в реда на трите подзадачи и ги тълкува. Всяко число
идва от изпълнение на програмата \code{sentinel\_bench} върху машината от
\tabref{tab:machine}, с командата

\begin{center}
\code{sentinel\_bench -{}-duration-ms 600 -{}-repeat 11}
\end{center}

\noindent
Пълният изход на това изпълнение, включително всички единадесет повторения на всеки ред, е
записан в хранилището като \code{docs/measurements/bench\_run.txt} и е приложен в
Приложение~\ref{sec:appendix}. Всяка обобщена стойност по-долу може да бъде проследена до
стойностите, от които е получена.

\subsection{Разчитане на ръкостискането и проверка на веригата}
\label{subsec:res_parser}

Резултатите от разчитането се дават в \tabref{tab:parser}. Докладваната стойност е медианата
от единадесет повторения, а разсейването е дадено и като междуквартилен размах, и като пълен
размах, по правилото от Раздел~\ref{subsec:validity}.

\begin{table}[H]
\caption{Пропускателна способност на разчитането и на проверката}
\label{tab:parser}
\centering
\fontsize{11}{13}\selectfont
\begin{tabular}{L{5.4cm} r r r r}
\toprule
Операция & ns за операция & операции/s & МКР & размах \\
\midrule
Разчитане на обмен по TLS 1.3, 2448 B & 7\,765,4 & 128\,776 & 11,5\,\% & 31,9\,\% \\
Разчитане на сертификат, 984 B & 8\,575,7 & 116\,608 & 8,0\,\% & 15,8\,\% \\
Изграждане и проверка на верига от 3 & 17\,330,4 & 57\,702 & 6,7\,\% & 13,7\,\% \\
\bottomrule
\end{tabular}
\end{table}

И трите реда са под своя праг от 25 на сто и се докладват. В други единици първият ред е
300,6 MiB/s върху обмена от 2448 байта, а вторият 109,4 MiB/s върху сертификата от 984 байта.

Проверката на верига от три сертификата струва 17,3 микросекунди, тоест два пъти повече от
разчитането на един сертификат. Разликата не е в разчитането, защото сертификатите се
разчитат веднъж и се подават разчетени: тя е в изграждането на пътя и в единадесетте
проверки върху него, като част от тях изграждат канонично представяне на отличителните имена
за сравнение. Числото е важно със своя порядък: при 57 хиляди проверки на верига в секунда на
едно ядро проверката на сертификати не е тясното място при никое натоварване, което една
машина може да поеме по TCP.

Присъдите по случаите от \tabref{tab:cases} съвпадат с очакваните за всичките тринадесет
случая. Изходът на \code{sentinel\_demo} завършва с реда „13 of 13 cases produced the verdict
they were built for“. Съответствието се проверява и автоматично от теста
\code{chain\_every\_generated\_case\_produces\_the\_verdict\_it\_was\_built\_for}, който е част
от набора и се изпълнява при всяка сборка. Случай, при който системата приема верига, която
трябва да отхвърли, няма; ако имаше, той щеше да бъде разгледан отделно, защото е дефект в
сигурността, а не разминаване в числата.

\tabref{tab:validator} показва състоянието на единадесетте проверки за изрядната верига.
Двата реда със състояние \code{SKIP} са същественото в тази таблица: те са проверките, които
не са проведени, и системата ги обявява вместо да ги преброи като успешни.

\begin{table}[H]
\caption{Състояние на всяка проверка при изрядна верига от три сертификата}
\label{tab:validator}
\centering
\fontsize{11}{13}\selectfont
\begin{tabular}{L{4.6cm} c L{7.0cm}}
\toprule
Проверка & Състояние & Забележка \\
\midrule
Свързване по име & \code{PASS} & три сертификата без прекъсване \\
Срок на валидност & \code{PASS} & всички покриват отправното време \\
Основни ограничения & \code{PASS} & всеки издател е удостоверител \\
Ограничение върху дължината & \code{PASS} & не е надхвърлено \\
Предназначение на ключа & \code{PASS} & краят: подпис и шифроване на ключ \\
Разширено предназначение & \code{PASS} & \code{serverAuth}, \code{clientAuth} \\
Съответствие на името & \code{PASS} & съвпада с \code{dNSName} \\
Ограничения върху имената & \code{SKIP} & никой издател по пътя не ограничава имена \\
Критични разширения & \code{PASS} & всички са разпознати \\
Отмяна & \code{PASS} & няма сериен номер в списъка \\
Подпис на издателя & \code{SKIP} & няма криптографска основа; двата подписа не са проверени \\
\bottomrule
\end{tabular}
\end{table}

\subsection{Цена на едно решение за допускане}
\label{subsec:res_decision}

\tabref{tab:decision} дава цената на едно решение вътре в процеса, без сокет по пътя. Това е
собствената цена на механизма, отделена от цената на връзката.

\begin{table}[H]
\caption{Цена на едно решение за допускане, измерена вътре в процеса}
\label{tab:decision}
\centering
\fontsize{11}{13}\selectfont
\begin{tabular}{L{5.6cm} r r r r}
\toprule
Конфигурация & ns за решение & МКР & размах & у клиента \\
\midrule
Без механизъм (база) & 8,7 & 40,7\,\% & 50,7\,\% & --- \\
Отчитане & 93,9 & 8,8\,\% & 15,1\,\% & --- \\
Отчитане + ограничаване & 94,9 & 9,4\,\% & 18,7\,\% & --- \\
Отчитане + бисквитка & 124,9 & 13,7\,\% & 31,9\,\% & --- \\
Отчитане + бисквитка + задача & 147,8 & 7,3\,\% & 49,0\,\% & 0,05 ms \\
Всички включени & 136,0 & 32,2\,\% & 54,2\,\% & 0,06 ms \\
\bottomrule
\end{tabular}
\end{table}

Първият ред надхвърля своя праг от 35 на сто и по правилото от Раздел~\ref{subsec:validity}
\textbf{не бива да бъде докладван като точна стойност}. Причината е ясна и не е в кода: при
базовата конфигурация решението се състои от увеличаване на брояч и връщане, тоест единици
наносекунди, а върху обща машина разсейването на планировчика е от същия порядък. Числото се
запазва тук с този надпис, вместо да бъде премахнато или загладено. Порядъкът му, единици
наносекунди, е достатъчен за извода, който се прави от него.

Останалите пет реда са в границите си. От тях се четат три величини:

\begin{itemize}[topsep=2pt,itemsep=2pt,leftmargin=*]
\item \textbf{Отчитането струва около 85 ns}, разликата между 8,7 и 93,9. Това е вписването и
      заличаването в хеш-таблицата с връзките, тоест цената на заделеното състояние.
\item \textbf{Ограничаването по източник струва около 1 ns над отчитането}, тоест 94,9 срещу
      93,9. Разликата е под разсейването на двата реда и от нея не бива да се вади заключение
      освен това, че цената е малка. Числото важи, когато ограничителят допуска.
\item \textbf{Бисквитката струва около 31 ns}, тоест 124,9 срещу 93,9. Това са две извиквания
      на SipHash-2-4, а при отказ на проверката и още едно за търпимия предходен прозорец.
\item \textbf{Проверката на изчислителната задача струва около 23 ns}, тоест 147,8 срещу
      124,9, което е едно хеширане.
\end{itemize}

Несиметричността на изчислителната задача е измерена пряко. При трудност от 12 бита търсенето
у клиента отнема средно \textbf{4279 хеширания} и \textbf{0,05 ms} на едно предизвикателство,
докато проверката у сървъра е едно хеширане, тоест около 23 ns. Отношението между двете е
малко над 2000. Ако трудността се вдигне с един бит, очакваният брой хеширания у клиента се
удвоява, а цената у сървъра остава същата: точно това е свойството, заради което механизмът
съществува, и то е причината политиката на плавно влошаване да вдига трудността, а не да
въвежда нова проверка.

Редът „Всички включени“ дава 136,0 ns, тоест по-малко от реда без ограничаване по източник.
Разликата е в рамките на разсейването на двата реда, чиито междуквартилни размахи са 32,2 и
7,3 на сто. От нея не се вади заключение. Отбелязва се, защото премълчаването на неудобно
подреден ред е точно начинът, по който една таблица става недостоверна.

\subsection{Поведение при лавинен трафик}
\label{subsec:res_admission}

\tabref{tab:admission} е основният резултат на проекта: цената на всеки механизъм под
действително натоварване, измерена спрямо базовата конфигурация. Всички редове са в границите
си по разсейване.

\begin{table}[H]
\caption{Цена на механизмите за допускане при лавинен трафик спрямо базата}
\label{tab:admission}
\centering
\fontsize{10}{12}\selectfont
\begin{adjustbox}{max width=\textwidth}
\begin{tabular}{L{3.9cm} r r r r r r}
\toprule
Конфигурация & допуснати/s & спрямо базата & p50, us & p95, us & върхова таблица & памет, B \\
\midrule
Без механизъм (база) & 4305 & 100,0\,\% & 454,2 & 575,9 & 0 & 0 \\
Отчитане & 3614 & 83,9\,\% & 524,9 & 762,8 & 287 & 7552 \\
Отчитане + ограничаване & 3450 & 80,1\,\% & 516,3 & 858,3 & 255 & 7912 \\
Отчитане + бисквитка & 1919 & 44,6\,\% & 987,8 & 1390,2 & 48 & 1504 \\
Отчитане + бисквитка + задача & 1671 & 38,8\,\% & 1094,2 & 1840,1 & 44 & 1408 \\
Всички включени & 1791 & 41,6\,\% & 1027,1 & 1692,8 & 43 & 1384 \\
\bottomrule
\end{tabular}
\end{adjustbox}
\end{table}

\subsubsection{Собствената цена на механизма е незначителна}
\label{subsubsec:res_negligible}

Съпоставката между \tabref{tab:decision} и \tabref{tab:admission} дава най-съществения извод
на измерването. Собствената цена на едно решение е от 9 до 148 наносекунди. Цената на една
допусната връзка под натоварване е от 454 до 1094 \emph{микро}секунди. Разликата е около
четири порядъка. Изчислението, което механизмът извършва, е невидимо на фона на работата по
установяването и затварянето на една TCP връзка.

Следователно намалението на пропускателната способност в \tabref{tab:admission} не идва от
изчислението. То идва от формата на протокола, което се вижда пряко в колоната с решенията за
секунда в необобщения изход: базата взема 12\,919 решения в секунда, а конфигурацията с
бисквитка взема 11\,532, тоест почти същото. Приемникът е еднакво зает и в двата случая. Това,
което се променя, е колко от тези решения завършват с допусната връзка.

\subsubsection{Бисквитката без състояние}
\label{subsubsec:res_cookie}

Бисквитката намалява допуснатите връзки до 44,6 на сто от базата и удвоява закъснението.
Причината е пряка: за да бъде допусната, една връзка изисква два обмена вместо един, тъй като
първият опит получава само предизвикателство и се затваря. Цената е една допълнителна TCP
връзка на допускане, а не изчисление.

Срещу тази цена се получава следното. Върховото запълване на таблицата пада от 287 записа на
48, тоест шест пъти, а заделената памет за състояние от 7552 на 1504 байта, тоест пет пъти,
при едно и също количество лавинен трафик. Четирите нишки с лавина в конфигурацията с
бисквитка не оставят нищо след себе си: всяка от техните връзки получава предизвикателство,
не отговаря и не води до заделяне. Същото свойство се проверява и пряко от теста
\code{harness\_a\_client\_that\_never\_answers\_the\_cookie\_leaves\_no\_state\_behind}, който
подава шестдесет неотговарящи и десет отговарящи клиента и изисква върхът на таблицата да не
надхвърли десет.

Това е сделката, която механизмът предлага, изразена в измерени числа: половината
пропускателна способност при нормални условия срещу шесткратно намаление на състоянието,
което лавината може да наложи. Дали сделката си струва, зависи от това дали ограничението е
пропускателната способност или паметта, и това е решение за конкретното разгръщане, не за
проекта.

\subsubsection{Ограничаване по източник}
\label{subsubsec:res_rate}

Ограничаването по източник дава 80,1 на сто от базата срещу 83,9 на сто за самото отчитане.
Разликата от 3,8 процентни пункта е в рамките на разсейването на двата реда, чиито
междуквартилни размахи са 11,6 и 13,2 на сто, и от нея не се вади заключение.

По-същественото наблюдение е качествено, а не числено. Върху локалния интерфейс честните
клиенти и лавинните клиенти делят един адрес на източника, така че ограничител, който
захапва, би задушил и двата вида еднакво. Това не е недостатък на стенда, а точното
проявление на известното ограничение на механизма: кофата с жетони по източник разделя
трафика само дотолкова, доколкото източниците са различими. Срещу разпределен източник тя не
работи, а срещу споделен адрес тя вреди. Че механизмът все пак действа, когато източниците са
различими, се проверява отделно от теста
\code{harness\_rate\_limiting\_refuses\_a\_single\_source\_that\_goes\_too\_fast}.

\subsubsection{Изчислителна задача и плавно влошаване}
\label{subsubsec:res_pow}

Добавянето на изчислителната задача върху бисквитката сваля допуснатите връзки от 44,6 на 38,8
на сто и вдига медианата на закъснението от 987,8 на 1094,2 микросекунди. Разликата от около
106 микросекунди се сравнява добре с измереното време за търсене у клиента от 0,05 ms, тоест
50 микросекунди, като остатъкът се пада на по-бавното обслужване при по-натоварен приемник.

Конфигурацията с всички механизми дава 41,6 на сто, тоест малко над конфигурацията само с
бисквитка и задача. Обяснението е в политиката на плавно влошаване: при това натоварване
запълването на таблицата остава ниско, режимът остава нормален и политиката не вдига
трудността. Тя не помага и не пречи при тези условия. Политиката се задейства при по-високо
запълване, което е проверено отделно от теста
\code{admission\_degradation\_raises\_the\_difficulty\_as\_the\_table\_fills}, но не се
проявява в тази конфигурация на опита. Това е ограничение на настоящото измерване и е записано
като такова, а не приписано на механизма.

\subsection{Отчитане на предварителните очакванията}
\label{subsec:res_hypotheses}

Всяко очакване от Раздел~\ref{subsec:hypotheses} се отчита изрично в \tabref{tab:hypotheses}.

\begin{table}[H]
\caption{Отчитане на предварителните очаквания спрямо измерените стойности}
\label{tab:hypotheses}
\centering
\fontsize{10}{12}\selectfont
\begin{tabular}{L{1.0cm} L{2.6cm} L{9.4cm}}
\toprule
Очак. & Изход & Измерване, което води до оценката \\
\midrule
О1 & потвърдено & върхова таблица от 287 на 48 записа, състояние от 7552 на 1504 байта, срещу допълнителен обмен и падане до 44,6\,\% от базата \\
О2 & потвърдено & 80,1\,\% срещу 83,9\,\% за самото отчитане, разлика в рамките на разсейването; върху споделен адрес механизмът не разделя честния от лавинния трафик \\
О3 & потвърдено & 4279 хеширания и 0,05 ms у клиента срещу едно хеширане и около 23 ns у сървъра, отношение над 2000 \\
О4 & потвърдено & 9 до 148 ns за решение срещу 454 до 1094 us за връзка, около четири порядъка разлика; решенията за секунда остават 12\,919 срещу 11\,532 \\
О5 & потвърдено & 7,765 us за пълен обмен, тоест 128\,776 обмена в секунда на едно ядро \\
\bottomrule
\end{tabular}
\end{table}

И петте очаквания се потвърждават. Това не е основание за задоволство: очакване, което се
потвърждава, е било или добре обосновано, или недостатъчно смело. Полезното от измерването не
е потвърждението на О1 до О5, а двете неща, които не бяха предвидени.

Първото е, че намалението на пропускателната способност при бисквитката се дължи изцяло на
допълнителния обмен, а не на изчислението, и това се вижда от почти непроменения брой решения
за секунда. Мерките за защита, които се обсъждат като „скъпи“, се оказват скъпи в друга
единица, отколкото се предполага: в брой обмени, не в процесорно време.

Второто е, че политиката на плавно влошаване не се задейства при тази конфигурация на опита,
защото бисквитката вече е свалила запълването на таблицата толкова, че режимът остава
нормален. Двата механизма са проектирани независимо, но при съвместна работа единият изважда
условието, при което другият има смисъл. Това е резултат, който се вижда само когато всички
механизми съществуват в една система и могат да бъдат включвани поотделно, тоест точно
заради решението от Раздел~\ref{subsec:admission}.
